\Askhsh{14884}{2}
\begin{ylh}{1}
    \item 3.2. 1ο Κριτήριο ισότητας τριγώνων
    \item 4.2. Τέμνουσα δύο ευθειών - Ευκλείδειο αίτημα
    \item 4.6. Άθροισμα γωνιών τριγώνου.
\end{ylh}
% ===================================================================
%  Αρχείο: 14884.tex (Fragment)
%  Αυτόματη μετατροπή μέσω doc2latex.py
% ===================================================================

ΘΕΜΑ 2

ΑΒΕ και ΑΔΖ.

Να αποδείξετε ότι:

\begin{alist}
\item Τα τρίγωνα ΑΕΖ και ΑΒΔ είναι ίσα. \monades{13}

\item Το τμήμα ΔΖ είναι παράλληλο στο ΒΕ. \monades{12}


\begin{center}
\begin{tikzpicture}[scale=0.8]
\tkzDefPoint(0,0){A}
\tkzDefPoint(4,0){B}
\tkzDefPointBy[rotation=center A angle 120](B) \tkzGetPoint{D_dir}
\tkzDefPointOnLine[pos=0.7](A,D_dir) \tkzGetPoint{D} % AD < AB for visual clarity

% Equilateral ABE exterior to ABD. A=0, B=4. E should be below
\tkzDefPointBy[rotation=center B angle -60](A) \tkzGetPoint{E}

% Equilateral ADZ exterior to ABD.
\tkzDefPointBy[rotation=center A angle -60](D) \tkzGetPoint{Z}
% Wait, "exterior". ABE is on side AB. Order is A-B-E. For ABE to be equilateral and exterior:
\tkzDefTriangle[equilateral](A,B) \tkzGetPoint{E_1}
\tkzDefTriangle[equilateral](B,A) \tkzGetPoint{E_2}
% D is above AB. So exterior ABE is E_2 (below AB).
\tkzDefPointBy[rotation=center A angle -60](B) \tkzGetPoint{E}

% For ADZ, D is at 120 degrees from B (y>0). Triangle ADZ should be exterior, so Z must be further left/down.
\tkzDefPointBy[rotation=center D angle -60](A) \tkzGetPoint{Z}

\draw[l3] (A)--(B)--(D)--cycle;
\draw[l3] (A)--(B)--(E)--cycle;
\draw[l3] (A)--(D)--(Z)--cycle;
\draw[l3, dashed] (A)--(E);
\draw[l3, dashed] (A)--(Z);
\draw[l3, dashed] (D)--(Z);
\draw[l3, dashed] (B)--(E);
% Wait, triangle matches the definition of ADZ equilateral, ABE equilateral.
\tkzDrawPoints(A,B,D,E,Z)
\tkzLabelPoint[below](A){$A$}
\tkzLabelPoint[right](B){$B$}
\tkzLabelPoint[above](D){$\Delta$}
\tkzLabelPoint[below right](E){$E$}
\tkzLabelPoint[left](Z){$Z$}
\end{tikzpicture}
\end{center}
\end{alist}
